\section{The enumerator method}
\label{sec:enumerator-method}

The minimum-distance event in \eqref{eq:bad-event} quantifies over all
nonzero messages.  We first count bad messages in expectation.  Exact
enumerators make the expectation computable without replacing the recursive
state process by a concentration inequality.

\subsection{Input--output weight enumerators}

Let $\mathcal D$ be a distribution on matrices in $\F_p^{a\times b}$.  Its
expected input--output weight enumerator is
\begin{equation}
 E^{\mathcal D}_{r,u}:=
 \E_{\M\gets\mathcal D}
 \left|\left\{\x\in\F_p^a:
   \wt(\x)=r,\ \wt(\x\M)=u
 \right\}\right|.
 \label{eq:io-enumerator}
\end{equation}
This is an expectation over matrices, not a count for one fixed matrix.

For a sampled generator $\G$, define
\begin{equation}
 Z_L(\G):=
 \left|\left\{\x\in\F_p^k\setminus\{0\}:
 \wt(\x\G)\leq L\right\}\right|.
 \label{eq:bad-message-count}
\end{equation}
The event $\Bad_L(\G)$ implies $Z_L(\G)\geq1$.  Therefore
\begin{equation}
 \Pr[\Bad_L(\G)]\leq \E[Z_L(\G)].
 \label{eq:first-moment}
\end{equation}
This is the only use of Markov's inequality in the common argument.  Our
improvement concerns how $\E[Z_L]$ is evaluated.

\subsection{Composition under exchangeability}

A monomial matrix has one nonzero entry in every row and column.  It permutes
coordinates and scales them by nonzero field elements.  A random vector whose
law is invariant under every monomial matrix is uniform on each Hamming shell
after conditioning on its weight.

\begin{lemma}[Exchangeable serial composition]
\label{lem:exchangeable-composition}
Let $\B\in\F_p^{k\times n}$ be sampled so that $\B\Pi$ and $\B$ have the same
distribution for every monomial matrix $\Pi\in\F_p^{n\times n}$.
Independently sample $\T\in\F_p^{n\times N}$.  Then
\begin{equation}
 E^{\B\T}_{r,h}=
 \sum_{u=0}^{n}
 \frac{E^{\B}_{r,u}E^{\T}_{u,h}}
      {\binom{n}{u}(p-1)^u}.
 \label{eq:exchangeable-composition}
\end{equation}
\end{lemma}

\begin{proof}
Fix a message $\x$ and condition on $\wt(\x\B)=u$.  Monomial invariance makes
$\x\B$ uniform on the $p$-ary weight-$u$ shell.  Averaging the independent
map $\T$ over that shell
gives the second enumerator divided by the shell size
$\binom{n}{u}(p-1)^u$.  Summing over $u$ and then over weight-$r$ messages
proves the identity.
\end{proof}

For binary matrices, coordinate-permutation symmetry alone suffices.  The
Bernoulli and global fixed-row ensembles satisfy the lemma.  A
region-stratified expander is exchangeable only within each region.  For that
ensemble we retain the recursive state at regional boundaries instead of
using \eqref{eq:exchangeable-composition} globally.

\subsection{Transfer matrices}

Suppose a recursive map has a finite state set $\mathcal S$.  Associate a
monomial $X^aY^b$ with a transition that consumes input weight $a$ and emits
output weight $b$.  Let $\mathbf T(X,Y)$ contain the transition probabilities
or expected transition counts.  If $\e_0$ selects the initial state, then
\begin{equation}
 [X^uY^h]\,\e_0^{\mathsf T}\mathbf T(X,Y)^n\one
 \label{eq:transfer-coefficient}
\end{equation}
is the expected mass of length-$n$ paths with input weight $u$ and output
weight $h$.  Matrix multiplication preserves the order of the recursive
process, while coefficient extraction forgets everything except the two
weights.

For a one-variable probability generating matrix $\mathbf M(Y)$ with
nonnegative coefficients, define
\[
 F(Y):=\e_0^{\mathsf T}\mathbf M(Y)^n\one.
\]
Every numerical marker $z\in(0,1]$ gives
\begin{equation}
 \sum_{h=0}^{L}[Y^h]F(Y)
 \leq z^{-L}F(z).
 \label{eq:positive-marker-tail}
\end{equation}
The left side uses coefficient notation schematically: the matrix on the
right is evaluated at the fixed marker $z$.  The inequality follows because
$z^h\geq z^L$ for $h\leq L$.  Optimizing $z$ improves the numerical bound,
but any fixed valid marker gives a sound bound.

\subsection{Covering all message weights}

A finite certificate must cover every $1\leq r\leq k$.  Evaluating selected
weights is not sufficient.  Our verifiers use three regimes:

\begin{enumerate}
  \item exact transfer or shell recurrences for small supports;
  \item one positive-marker bound for each contiguous support interval;
  \item a point-mass or Hamming-ball bound for dense supports.
\end{enumerate}

Every certificate records an explicit partition of $[1,k]$.  Its verifier
rejects overlapping intervals, gaps, invalid markers, or a final sum above the
claimed value.  Section~\ref{sec:certificates} describes the trusted
arithmetic.

The following elementary lemma isolates the common block calculation.  The
theorem-specific analysis supplies the envelope and proves its shape.  The
verifier only evaluates the stated endpoints and checks the partition.

\begin{lemma}[Certified support partition]
\label{lem:certified-support-partition}
Let $f_r\geq0$ be the first-moment contribution of support $r$, and let
$\mathcal I$ be a disjoint partition of $\{1,\ldots,k\}$ into integer
intervals.  For each $I=[a,b]\in\mathcal I$, suppose a fixed valid marker
defines an upper envelope $F_I:[a,b]\to\mathbb R_{\geq0}$ with
$f_r\leq F_I(r)$ for every integer $r\in I$.  Define
\begin{equation}
 M_I:=
 \begin{cases}
  F_I(b),&F_I\text{ is nondecreasing},\\
  F_I(a),&F_I\text{ is nonincreasing},\\
  F_I(a)+F_I(b),&F_I>0\text{ and }\ln F_I\text{ is convex}.
 \end{cases}
 \label{eq:certified-block-endpoint}
\end{equation}
Any applicable line gives
\begin{equation}
 \sum_{r=1}^{k}f_r
 \leq\sum_{I=[a,b]\in\mathcal I}(b-a+1)M_I.
 \label{eq:certified-partition-sum}
\end{equation}
The same conclusion holds when a theorem supplies a direct upper bound
$B_I\geq\sum_{r\in I}f_r$, with $B_I$ replacing $(b-a+1)M_I$.
\end{lemma}

\begin{proof}
Monotonicity gives the first two endpoint bounds.  If $\ln F_I$ is convex,
its maximum on the real interval $[a,b]$ occurs at an endpoint.  Thus
$F_I(r)\leq\max\{F_I(a),F_I(b)\}\leq F_I(a)+F_I(b)$.
Summing the pointwise bound over each interval proves
\eqref{eq:certified-partition-sum}.  The direct-bound case follows by
summing the inequalities $\sum_{r\in I}f_r\leq B_I$.
\end{proof}

An interval verifier implements this lemma with outward-rounded enclosures.
It checks that the intervals partition $[1,k]$, that each marker is in its
proved domain, and that the upper endpoint of every computed $M_I$ or $B_I$
is finite.  Symbolic monotonicity and convexity remain mathematical claims;
the verifier invokes them only through the named bounds identified in
Section~\ref{sec:certificates}.
